\chapter*{Как читать это пособие}
\addcontentsline{toc}{chapter}{Как читать это пособие}

Это пособие учит проектировать цифровые устройства. Оно рассчитано на тех,
кто только начинает знакомство с цифровой электроникой: достаточно школьной
математики и желания разобраться, как «думают» компьютеры, контроллеры
дронов и любые другие цифровые устройства.

\section*{Идея книги: проектируй, запусти, наблюдай}

Цифровую схему нельзя выучить, только читая о ней. Поэтому в каждой главе
вы повторяете один и тот же цикл:

\begin{enumerate}
  \item \textbf{Проектируй.} Соберите схему в Конструкторе Logra или
    опишите её в Коде.
  \item \textbf{Запусти.} Задайте входы и запустите симуляцию.
  \item \textbf{Наблюдай.} Сравните диаграмму с таблицей из книги, посмотрите схему
    на вкладке RTL и объясните, почему результат получился именно таким.
\end{enumerate}

Logra~--- это среда, в которой вы работаете с первой страницы. Она состоит
из вкладок Обучение, Код, Конструктор, Выводы, RTL, Диаграммы, Лаборатория,
Gowin, Кристалл и Подключение. В боковой панели они собраны в группы
«Узнать», «Создать», «Проверить» и «Настоящая плата». Начинать работу с ней
можно сразу: первые главы опираются только на Конструктор, Диаграммы и RTL.

\section*{Как устроена глава}

Каждая глава~--- небольшой «спринт». Она начинается с мостика: что мы уже умеем и
какой вопрос остался. Затем идёт врезка «Соберите сначала» с заданием в
Logra, которое стоит выполнить до чтения теории: так вы встретите
теорию уже с вопросом. Каждый раздел заканчивается заданием
\textbf{«Сделай сам»}: что построить в Logra, какой результат ожидать и как
его проверить. В конце главы идут типичные ошибки, вопросы и задачи и мостик к
следующей теме.

\begin{infobox}[Правило главы]
Не переходите к следующему разделу, пока диаграмма не сходится с таблицей.
Если сигналы на вкладке Диаграммы расходятся с таблицей истинности, то
ошибка либо в схеме, либо в таблице. Найдите её сейчас: дальше она только
накопится.
\end{infobox}

\subsection*{Ступени сложности}

Задания и задачи помечены ступенями:
\stars{1}~--- \emph{повтори}: сделать по образцу и сверить результат;
\stars{2}~--- \emph{измени}: изменить готовую схему или найти в ней ошибку;
\stars{3}~--- \emph{проект}: решить задачу самостоятельно и принять результат по
заданному критерию. Если трудно, вернитесь на ступень ниже, это нормально.

\section*{Сквозной проект «Охранная панель»}

Через всю книгу вы строите одно устройство: охранную панель с датчиками двери и окна,
ночным режимом, ключом охраны, кодовым замком, сиреной и индикатором. Каждая
глава добавляет к ней новую часть (датчик как бит, формулу тревоги,
голосование датчиков, выбор зоны, порог шума, код замка, запоминание тревоги,
счётчик времени, автомат охраны), и панель всегда остаётся рабочей. К концу
теоретической линии у вас готовое устройство, а в главах про плату и проекты оно оказывается
на ПЛИС.
Этапы собраны в приложении~\ref{app:logramap}.

\section*{Как устроена книга}

Книга одна, сквозная, без деления на части. Теория, Verilog и работа в Logra
идут вместе: вы проектируете, запускаете и наблюдаете на каждом шаге.

\begin{description}
  \item[Старт.] Что такое ПЛИС и чем она отличается от процессора; первый
    запуск в Logra: вкладки, цикл «проектируй, запусти, наблюдай».
  \item[Спринты по темам.] Сигналы; системы счисления и коды; логические
    элементы и булева алгебра; язык Verilog; комбинационные схемы;
    минимизация (карты Карно, метод Квайна---Мак-Класки); дешифраторы и
    шифраторы; мультиплексоры; арифметические узлы; триггеры; регистры и
    счётчики; синхронные схемы и конечные автоматы; проектирование автоматов и
    очереди FIFO. В каждой главе теория, код на Verilog и схема в Logra
    рассматриваются вместе.
  \item[Плата, маршрут, физика, аналог.] Что внутри ПЛИС (LUT, триггеры,
    трассировка); плата Sipeed Tang Nano 20K; путь от текста до прошивки;
    электрические основы КМОП и шин; ЦАП, АЦП и ШИМ; мост от схемы к Verilog.
  \item[Проекты и системы.] Семь законченных проектов; симуляция и
    тестбенчи; тактирование, временной анализ и пересечение тактовых доменов;
    интерфейсы UART, SPI и I$^2$C; память SDRAM; собственный процессор,
    который запускает программу; задания для самостоятельной работы.
\end{description}

\begin{figure}[H]
\centering
\begin{tikzpicture}[node distance=5mm and 4mm, t/.style={block, minimum width=28mm, minimum height=15mm, font=\sffamily\footnotesize}]
  \node[t] (a) {Старт:\\ПЛИС и Logra};
  \node[t, right=of a] (b) {Сигналы,\\элементы};
  \node[t, right=of b] (c) {Verilog,\\комбинационные\\узлы};
  \node[t, right=of c] (d) {Триггеры,\\регистры};
  \node[t, right=of d] (e) {Автоматы,\\FSM};
  \node[hblock, minimum width=28mm, minimum height=15mm, font=\sffamily\footnotesize, below=9mm of e] (f) {Плата,\\маршрут};
  \node[hblock, minimum width=28mm, minimum height=15mm, font=\sffamily\footnotesize, left=of f] (g) {Физика,\\аналог};
  \node[hblock, minimum width=28mm, minimum height=15mm, font=\sffamily\footnotesize, left=of g] (h) {Мост к\\Verilog};
  \node[oblock, minimum width=28mm, minimum height=15mm, font=\sffamily\footnotesize, left=of h] (i) {Проекты,\\UART, такты};
  \node[oblock, minimum width=28mm, minimum height=15mm, font=\sffamily\footnotesize, left=of i] (j) {SDRAM, CPU,\\задания};
  \foreach \u/\v in {a/b, b/c, c/d, d/e, e/f, f/g, g/h, h/i, i/j} \draw[arr] (\u) -- (\v);
\end{tikzpicture}
\caption{Путь, который мы пройдём}
\end{figure}

Если вы уже знаете логику, триггеры и автоматы, можно сразу переходить к главам
про плату и проекты (начиная с главы~\ref{ch:inside}) и возвращаться к
спринтам как к справочнику: в них есть ссылки на нужные главы. Соответствие глав урокам и примерам Logra
сведено в приложении~\ref{app:logramap}.

Все исходные тексты примеров лежат рядом с книгой, в папке
\nolinkurl{textbook/code}. Их можно открыть, скопировать и
сразу загрузить в плату.

\begin{infobox}[Обозначения]
\begin{itemize}
  \item Синие блоки «Справка» содержат полезные подробности.
  \item Оранжевые блоки «Внимание!» предупреждают о типичных ошибках.
  \item Фиолетовые блоки «Главная мысль» подводят итог раздела.
  \item Зелёные блоки «Задание» и «Сделай сам» предлагают что-то сделать
    самостоятельно в Logra.
  \item Голубые блоки «Соберите сначала» дают задание до теории.
  \item Курсивом «Заглянем вперёд» отмечены редкие взгляды на поздние главы.
\end{itemize}
\end{infobox}
